\documentclass[twocolumn]{aastex631}
\begin{document}
\title{A residual reionization ionizing-photon-budget shortfall for star-forming galaxies at z\~{}6 (crisis systematic corner)}
\author{NebulaMind Lab (autonomous pipeline)}
\affiliation{NebulaMind Lab, \url{https://lab.nebulamind.net}}
\begin{abstract}
Reionization ionizing-photon-budget reconciliation at z\~{}6: star-forming galaxies require f\_esc=0.180 (+0.143/-0.079) to close the budget (Madau-Dickinson SFRD, log xi\_ion=25.5+/-0.15, clumping C=2-5, JWST-SFRD tail), versus indirect-proxy-inferred f\_esc=0.050 (+0.075/-0.030) from LzLCS O32/beta calibrations. Median delta(required-inferred)=+0.119 dex-frac (16-84\%: +0.020 to +0.264); 88\% of the systematic MC shows a shortfall. Conclusion: a genuine SHORTFALL remains, and the sign holds under both O32 and beta calibrations. The result is bounded by the xi\_ion x clumping x proxy-calibration systematic, not by statistics. Generated autonomously from public data (jwst) via the NebulaMind Lab runner: a bounded, reproducible, descriptive study.
\end{abstract}
\section{Introduction}
Recent studies have highlighted a potential crisis in understanding the reionization process, with concerns that current models may not account for the necessary ionizing photons to drive this cosmic event [Muñoz2024]. This issue is further complicated by the need for increased demands on ionizing sources, as discussed in Davies et al. (2021). To address these challenges, researchers have explored various approaches, including excursion set reionization models [Park2022] and assessments of the galaxy ionizing photon budget at high redshifts [Duncan2015]. However, a comprehensive understanding of the reionization process remains elusive.
\section{Data and method}
In this work, we adopt a literature-anchored approach to calculate the reionization-photon-budget. We utilize the cosmic star formation rate density (SFRD) from Madau \& Dickinson (2014), along with published values for xi\_ion and O32/beta f\_esc proxy calibrations [Chisholm+22, Flury+22; Simmonds+24]. Our method focuses on reconciling systematic uncertainties in the ionizing-photon-budget using these literature values.
\section{Result}
Our analysis reveals that star-forming galaxies require an escape fraction of f\_esc=0.180 (+0.143/-0.079) to close the reionization photon budget at z\~{}6, assuming a Madau-Dickinson SFRD, log xi\_ion=25.5+/-0.15, and clumping factor C=2-5. In contrast, indirect-proxy-inferred f\_esc values from LzLCS O32/beta calibrations yield f\_esc=0.050 (+0.075/-0.030). This results in a median delta of +0.119 dex-frac (16-84\%: +0.020 to +0.264), with 88\% of systematic Monte Carlo simulations showing a shortfall.
\begin{figure}\centering\includegraphics[width=\columnwidth]{result.png}\caption{A residual reionization ionizing-photon-budget shortfall for star-forming galaxies at z\~{}6 (crisis systematic corner)}\end{figure}
\section{Caveats}
It is essential to acknowledge the limitations of our approach, which relies on automated single-selection and uncalibrated measurements. The accuracy of our results depends heavily on the assumptions made in the literature values we adopt, particularly regarding xi\_ion and O32/beta f\_esc proxy calibrations. Furthermore, our analysis does not account for potential variations in these parameters across different galaxy populations or redshift ranges. These factors introduce uncertainties that must be considered when interpreting our findings.
\section*{References}
\footnotesize
[Muoz2024] Muñoz et al. (2024). Reionization after JWST: a photon budget crisis?. bibcode 2024MNRAS.535L..37M \par [Davies2021] Davies et al. (2021). The Predicament of Absorption-dominated Reionization: Increased Demands on Ionizing Sources. bibcode 2021ApJ...918L..35D \par [Park2022] Park et al. (2022). Calibrating excursion set reionization models to approximately conserve ionizing photons. bibcode 2022MNRAS.517..192P \par [Duncan2015] Duncan et al. (2015). Powering reionization: assessing the galaxy ionizing photon budget at z < 10. bibcode 2015MNRAS.451.2030D \par [Madau2017] Madau (2017). Cosmic Reionization after Planck and before JWST: An Analytic Approach. bibcode 2017ApJ...851...50M

\end{document}
